\section{MIT 8.323 Relativistic Quantum Field Theory 1}

\subsection{Lecture 1: Classical Field Theories and Principle of Locality}
The first part goes over trivial procceses in field theory and general mathematics. The main note on this section would be the principle of Locality:
\[
L = \int d^3 x \mathcal{L} \left(\phi_a , \partial_t \phi_a , \partial_i \phi_a  \right)
  \]

  Where $\phi_a(\vec{x}, t)$: general fields.
$\mathcal{L}$ is a function of $\phi$ and its derivatives. $\mathcal{L}$ is called the Lagrangian density.

For this form it constrains all dynamics so that all actions must be local. It is used by a single integral to enforce this. Another constraint is one will only use the first derivative in time. This is because including a second order time derivative will make an equation of motion with an initial condition with an acceleration which includes arbitrary complications without need. This rule is less important and fundamental but is used to make things easier. Lastly, in this class, being that it is relativistic one won't need additional spacial derivatives, but in classical models this in not the case.

Canonical momentum density is exactly what one would expect.
\[
\Pi_a = \frac{\partial \mathcal{L}}{\partial \phi_a}
  \]
  Hamiltonian is defined the same way.

  One brief note on notation, in this x is simply a label for a degree of freedom, there can be infinite such degrees of freedom. Also, the teacher is using Einstein's notation that repeated indices will be summed.


\subsubsection{Equations of motion}
He gives an example through:
\[
  \delta s = 0
\]
\[
S = \int dt L
  \]
\[
0 = \delta s = \int d^4 x \left[ \frac{\partial \mathcal{L}}{\partial \phi_a} \delta_a + \frac{\partial \mathcal{L}}{\partial(\partial_\mu \phi_a)} \delta(\partial_\mu \phi_a\right]
  \]
  One can simplify this through integration by parts

  As for boundary conditions. One will choose $\delta \phi_a $ so that boundary conditions will vanish. Thus one will not need to worry about that.

This thus implies that $ \frac{\partial \mathcal{L}}{\partial \phi_a} - \partial_\mu \left( \frac{\partial \mathcal{L}}{\partial(\partial_\mu \phi_a)} \right) = 0$.

From now on all things will be translationally invariant and Lorentz invariant.

He briefly goes over the simplest form of a quantum field, a scalar field.(Pion or higgs)

It follows a very similar model:
\[
S = \int d^4 x \left[ -1/2, \partial_\mu \phi, \partial^\mu \phi - V(\phi) \right]
  \]

  Transnational symmetry means that all parameters of $V(\phi)$ must be constant. V is the same as the Lagrangian density

Using this equation one can once again find the Hamiltonian Density and thus Hamiltonian. To give a simple equation of motion.

He goes through a famous Klein-Gordan:
\[
\partial^2 \phi - m^2 \phi = 0
  \]
\subsection{Lecture 2: Symmetries and Conservation Laws}

I forgot to save my notes last time, so those are gone. It was generally pretty simple going over Noether theorem and defining types of symmetries at the beginning.

Some how it deleted again... oh well, moving on.

\subsection{Lecture 3: Why Quantum Field Theory}

Non-Relativistic Schrodinger equation

\[
\partial^2 \psi - m^2 \psi = 0
  \]
  Simplest scalar field theory
\[
S= - \frac 12 \int dx (\partial_\mu \phi \partial^\nu \phi + m^2 \phi^2)
  \]
  Implies
\[
\partial^2 \phi -m^2 \phi = 0
  \]

  This changes $\phi(t, x)$ where x becomes an eigenvalue rather than a coordinate.

In relativistic quantum mechanics, there are several challenges. For one, no way to define a positive probability density. Also, there exists negative energy states. You must describe a fixed number of particles without changing. As equations change as wavefunctions change. This troubles with relativity as energy can create particles, adding additional difficulty. Finally, there is a fundamental asymmetry between the parameter t and the eigenvalues of x.

Therefore, relativistic quantum mechanics one cannot give a fundamental description.(at best an approximate, no particle creation or annihilation) The solution of this is quantum field theory, simply the quantization of classical field theory.(using t and x as parameters not eigenvalues) Field theory can also arise as a limit of discrite systems.

In classical field theory, one takes the originally discrite Lagrangian and defines $\phi(x_n,t)$ to be the discrite parts. From here a simple Taylor expanse can be used to change all of the interaction terms to in the form of the new continuum. Finally one may derive the field-theoretic Euler-Lagrange equation.

Quantization of harmonic oscillator in Heisenberg picture.

\[
L=\frac 1 2 \dot{x} - \frac{\omega^2}{2}x^2
  \]
\[
p=\dot{x}
  \]
\[
H=\frac{p^2}{2}+\frac 1 2 \omega^2x^2
  \]
\[
\ddot{x}+\omega^2x=0
  \]
  The quantum solution is exactly the same, exept x becomes an operator:
\[
\hat{x}(t)= \hat{A}\cos{\omega t} + \hat{B} \sin{\omega t}
  \]
  From here one can define the hilbert space to define through $ [\hat{a}, \hat{a^+} ] i$
\[
|h>= \frac{1}{\sqrt{n!}}(\hat{a^+})^n | 0>
  \]

  One can then use this to quantize a field theory. First, find the most general solution of the EOM. Second, promote the integration constants to constant operators to give full time evolution. Impose canonical quantization conditions$\{q,p \}=1$ thus, $ [\hat{q}, \hat{p} ]= ih$. Finally, use this to develop Hilbert Space.

\subsection{Lecture 4: Canonical Quantization of a Free Scalar Field Theory}
